\chapter{Уся машина}
\chaptermark{Уся машина}
\label{sec:synthesis}

\section{Що ми зібрали}

Ми брали з таблиці~\ref{tab:types} по одному типу нейронів і питали, що він додає до обчислювальної машини. Щоразу правила були ті самі: лише те, що буває в живому організмі, і жодного спільного тактового сигналу. Тепер зберемо частини в одну машину й подивимося, як вони працюють разом.

Картина з розділу~\ref{sec:neurontypes} підтвердилася і стала точнішою. Величезна адресна мережа \nt{GLUT}-нейронів несе дані. \nt{GABA}-нейрони керують потоком цих даних. А над мережею висять чотири широкомовні канали, кожен зі своєю ручкою: \nt{DOPA} каже, які зміни ваг залишити, \nt{SERO} тримає загальний рівень і, за гіпотезою, горизонт планування, \nt{NORA} вибирає між пошуком і рішенням, \nt{ACET} --- між записом і читанням (рис.~\ref{fig:machine}).

\begin{figure}[t]
\centering
\begin{tikzpicture}[>=Stealth,
  blk/.style={draw,very thick,rounded corners=3pt,align=center,font=\small},
  reg/.style={draw,very thick,double,double distance=1.2pt,circle,minimum size=1.25cm,inner sep=0pt,font=\scriptsize,fill=orange!15},
  bc/.style={->,thick,densely dashed,orange!85!black}]
% sensors, bus, actions
\node[blk,fill=green!8,text width=1.7cm] (S) at (0.9,0) {датчики\\\scriptsize увімк./вимк.};
\draw[very thick,rounded corners=4pt,fill=blue!5] (2.6,-1.35) rectangle (9.0,1.35);
\node[font=\small] at (5.8,0.95) {шина \nt{GLUT}: дані};
\node[font=\scriptsize,align=center] at (4.3,0.25) {збіги, затримки,\\логіка (розд.~\ref{sec:glutcomputer})};
\node[font=\scriptsize,align=center] at (7.4,0.25) {пам'ять, код\\(розд.~\ref{sec:code}, \ref{sec:acet})};
\draw[thick,red!70!black,fill=red!8,rounded corners=2pt] (2.8,-1.15) rectangle (8.8,-0.55);
\node[font=\scriptsize,red!60!black] at (5.8,-0.85) {\nt{GABA} (розд.~\ref{sec:gaba}): заборона, вибір, ділення, ритм};
\node[blk,fill=blue!5,text width=2.0cm] (A) at (10.9,0) {вибір\\дії};
\draw[->,very thick] (S) -- (2.6,0);
\draw[->,very thick] (9.0,0) -- (A);
\draw[->,very thick] (A) -- (12.6,0) node[right,font=\small] {м'язи};
\pic[scale=1.1] at (13.2,-0.5) {spring};
\pic[scale=1.15] at (0.4,0.85) {eyepic};
\pic[scale=1.05] at (1.35,0.85) {speaker};
% registers
\node[reg] (D) at (3.4,3.2) {\nt{DOPA}};
\node[reg] (C) at (5.6,3.2) {\nt{SERO}};
\node[reg] (N) at (7.8,3.2) {\nt{NORA}};
\node[reg] (H) at (10.0,3.2) {\nt{ACET}};
\node[font=\scriptsize,align=center,above=1pt] at (D.north) {помилка $\delta$};
\node[font=\scriptsize,align=center,above=1pt] at (C.north) {рівень, $\tau_\gamma$};
\node[font=\scriptsize,align=center,above=1pt] at (N.north) {підсилення $g$};
\node[font=\scriptsize,align=center,above=1pt] at (H.north) {запис $a$};
\draw[bc] (D.south) -- (3.4,1.35);
\draw[bc] (C.south) -- (5.6,1.35);
\draw[bc] (N.south) -- (7.8,1.35);
\draw[bc] (H.south) -- (10.0,2.1) -- (8.8,1.35);
\draw[bc] (H.south east) to[bend left=20] (A.north east);
\draw[->,thick] (2.85,1.35) -- (2.85,2.75) -- (D.west) node[pos=0.25,left,font=\scriptsize] {$V$};
\draw[->,thick,green!45!black] (0.4,3.2) node[left,black,font=\small] {винагорода $r$} -- (2.2,3.2) -- (D.west);
\pic[scale=0.9] at (-0.45,3.7) {juice};
\draw[->,thick,densely dotted,gray!50!black] ([yshift=0.45cm]D.north) to[bend left=18] node[above,font=\scriptsize,black] {несподіванка} ([yshift=0.45cm]N.north);
\end{tikzpicture}
\caption{Машина цілком. Шина \nt{GLUT} несе дані від датчиків до вибору дії; \nt{GABA} керує потоком усередині неї. Подвійні кружки --- широкомовні канали; штрихові стрілки --- розсилання їхніх сигналів усій мережі, у кожного своя ручка. \nt{DOPA} отримує винагороду $r$ і очікування $V$ від критика всередині шини (розд.~\ref{sec:dofa}); точкова стрілка --- гіпотеза, що про несподіванку \nt{NORA} дізнається з незвично великих помилок (розд.~\ref{sec:nora}).}
\label{fig:machine}
\end{figure}

\section{Одна формула для всіх ручок}

Ручки можна зібрати в одну схематичну формулу. Це не нова модель, а зведення моделей розділів~\ref{sec:glutcomputer}--\ref{sec:acet}: кожен символ узято зі свого розділу.
\begin{empheq}[box=\keybox]{equation}
\begin{aligned}
\tau\,\frac{\dd r_i}{\dd t} \;&=\; -r_i + F\Bigl(g\,\Bigl[\tfrac{1+a}{2}\,x_i + (1-a)\sum_j W_{ij}\,r_j - \sum_k M_{ik}\,q_k - \theta\Bigr]\Bigr),\\
\frac{\dd W_{ij}}{\dd t} \;&=\; \eta\,a\,\delta(t)\,e_{ij}(t),
\qquad \tau_e\,\frac{\dd e_{ij}}{\dd t} = -e_{ij} + r_i\,r_j,
\qquad \delta = r + \frac{\dd V}{\dd t} - \frac{V}{\tau_\gamma} .
\end{aligned}
\label{eq:machine}
\end{empheq}
Тут $r_i$ --- частоти \nt{GLUT}-нейронів, $x_i$ --- вхід від датчиків, $W_{ij}\ge0$ --- їхні взаємні ваги (розділ~\ref{sec:glutcomputer}); $q_k$ --- частоти \nt{GABA}-нейронів, $M_{ik}\ge0$ --- сила їхнього гальмування (розділ~\ref{sec:gaba}); $e_{ij}$ --- слід, $\delta$ --- помилка \nt{DOPA} (розділ~\ref{sec:dofa}); $\tau_\gamma$ --- горизонт, який, за гіпотезою, задає \nt{SERO}; $g$ --- підсилення \nt{NORA} (розділ~\ref{sec:nora}); $a$ --- рівень \nt{ACET} (розділ~\ref{sec:acet}).

Погляньте, чого в~\eqref{eq:machine} немає. Немає такту: усе записано диференціальними рівняннями, і кожен нейрон змінюється сам. Немає окремої пам'яті: пам'ятають ті самі $W_{ij}$, якими йде обчислення, і та сама активність $r_i$. Немає персональних поправок для переважної більшості синапсів: їхнє навчання --- добуток місцевого сліду й одного спільного числа; окремий провід помилки до кожного виходу є лише в колі, що навчає рухів (розділ~\ref{sec:motorlearning}). І керувальних сигналів лише чотири види: $\delta$, $\tau_\gamma$, $g$, $a$ на вісімдесят мільярдів нейронів. Кожен із них насправді не одне число, а невеликий джгут паралельних ліній (твердження~\ref{prop:bundle}), але ліній у джгуті лічені одиниці.

\begin{keyblock}
\begin{proposition}[Архітектура]
\label{prop:architecture}
Машину мозку, наскільки ми її зараз розуміємо, можна описати трьома шарами. Дані йдуть адресною шиною \nt{GLUT}. Потік даних спрямовують, обмежують і нормують адресні \nt{GABA}-нейрони. Режим роботи задають кілька широкомовних каналів, кожен --- невеликий джгут паралельних ліній, спільних для великих ділянок мережі. Перші два шари встановлено твердо; розподіл ролей між широкомовними каналами --- здебільшого гіпотеза.
\end{proposition}
\end{keyblock}

\section{Усі типи в одній таблиці}

Таблиця~\ref{tab:synthesis} зводить усі типи нейронів книги. Чотири типи, яким не дістався власний розділ, посіли в ній по рядку; двоє з них знайшли роль у розділах про виходи машини: \nt{GLYC} --- у петлях спинного мозку (розділ~\ref{sec:muscles}), \nt{ADRE} --- у хімічному виході (розділ~\ref{sec:chemoutput}). Роль решти двох в обчисленні або проста, або невідома. Речовини, що не задають тип нейрона, зібрано в додатку~\ref{sec:othermediators}.

\begin{table}[t]
\centering
\footnotesize
\setlength{\tabcolsep}{3pt}
\renewcommand{\arraystretch}{1.15}
\begin{tabular}{>{\raggedright}p{1.3cm}>{\raggedright}p{1.0cm}>{\raggedright}p{3.9cm}>{\raggedright}p{1.5cm}>{\raggedright}p{1.8cm}>{\raggedright\arraybackslash}p{3.8cm}}
\toprule
Тип & Розділ & Що обчислює & Час & Шина & Що відомо \\
\midrule
\nt{GLUT} & \ref{sec:glutcomputer}, \ref{sec:code}, \ref{sec:sensors}, \ref{sec:motorlearning}, \ref{sec:dendrites} & дані: збіги, затримки, логіка, пам'ять, послідовності & мс & адресна & встановлено \\
\nt{GABA} & \ref{sec:gaba}, \ref{sec:code}, \ref{sec:pain}, \ref{sec:motorlearning} & керування потоком: заборона, вибір, ділення, стійкість, ритм; ворота болю & мс & адресна & встановлено; ділення частоти --- разом із фоновим шумом \\
\nt{SERO} & \ref{sec:serocomputer}, \ref{sec:pain}, \ref{sec:sleep} & регулятор рівня, згладжування; горизонт $\tau_\gamma$; гучність болю; режими сну & секунди & об'єм & самогальмування виміряно; горизонт --- гіпотеза \\
\nt{DOPA} & \ref{sec:dofa}, \ref{sec:dofaalt} & помилка передбачення $\delta$; повільний рівень --- ціна часу & $0.1$\,с і хвилини & широкомовна & $\delta$ виміряно; розширення --- розд.~\ref{sec:dofaalt} \\
\nt{NORA} & \ref{sec:nora}, \ref{sec:pain}, \ref{sec:chemoutput}, \ref{sec:sleep} & підсилення $g$: шукати чи вирішувати; несподіванка; «газ» для органів & секунди & широкомовна & зростання відношення сигнал/шум виміряно; роль у пошуку --- гіпотеза \\
\nt{ACET} & \ref{sec:acet}, \ref{sec:muscles}, \ref{sec:chemoutput}, \ref{sec:sleep} & запис чи читання; швидкість навчання; вихід на м'яз; «гальмо» для органів & секунди & обидві & три ефекти виміряно; перемикання режимів --- гіпотеза \\
\nt{GLYC} & \ref{sec:muscles}, \ref{sec:pain} & від'ємна вага в спинному мозку й стовбурі: заборона м'яза-антагоніста & мс & адресна & встановлено \\
\nt{HIST} & --- & глобальний сигнал «не спи» & повільно & широкомовна & роль в обчисленні не вивчено \\
\nt{ADRE} & \ref{sec:chemoutput} & «газ» для всього тіла через кров; у мозку невідомо & повільно & широкомовна & поза мозком встановлено; роль у мозку неясна \\
\nt{ASPA} & --- & слабка додатна вага & мс & адресна & роль спірна \\
\bottomrule
\end{tabular}
\caption{Усі типи нейронів книги: що обчислює кожен і наскільки це відомо.}
\label{tab:synthesis}
\end{table}

\section{Як ручки працюють разом}

Досі кожен розділ повертав одну ручку. Простежмо тепер одну подію через усю машину (рис.~\ref{fig:timeline}). Тварина давно вивчила, що дія $A$ приносить сік. Якоїсь миті світ змінюється: $A$ більше не приносить соку, а приносить його дія $B$, яку тварина майже не пробує.

\emph{Помилка.} Сік після $A$ не прийшов, і \nt{DOPA}-нейрони відповідають провалом, $\delta<0$ за формулою помилки~\eqref{eq:ctd}. Синапси, які щойно вибрали $A$, несуть слід і слабшають.

\emph{Несподіванка.} Один провал --- звичайна випадковість. Але провали повторюються, і помилки стають більшими за звичайні. За гіпотезою розділу~\ref{sec:nora}, це і є сигнал \nt{NORA}: світ змінився. Підсилення $g$ падає, вибір стає м'яким, і шум частіше приводить до спроби $B$.

\emph{Запис.} За гіпотезою розділу~\ref{sec:acet}, після зміни \nt{ACET} зростає й переводить мережу в режим запису: внутрішні зв'язки, що зберігають стару звичку, послаблено, вхід від датчиків підсилено, ваги змінюються легко.

\emph{Нова звичка.} Спроба $B$ приносить сік, якого ніхто не чекав: сплеск, $\delta>0$, і синапси, що вибрали $B$, посилюються. Кілька таких спроб --- і очікування $V$ наздоганяє новий світ, помилки знову малі.

\emph{Повернення.} Несподіванка минула, \nt{NORA} повертає високе підсилення, вибір знову жорсткий; \nt{ACET} спадає, мережа в режимі читання користується вивченим. Увесь цей час \nt{SERO}, за гіпотезою, задає горизонт $\tau_\gamma$: наскільки далекий сік ще варто чекати.

\begin{figure}[t]
\centering
\begin{tikzpicture}[>=Stealth,x=1.1cm,y=0.95cm]
\foreach \y/\lab in {4/{$\delta$ (\nt{DOPA})},3/{$g$ (\nt{NORA})},2/{$a$ (\nt{ACET})},1/{вибір $B$}}{
  \draw[gray!60!black] (0,\y-0.35) -- (10,\y-0.35);
  \node[left,font=\scriptsize] at (0,\y) {\lab};}
\draw[densely dashed,gray!60!black] (2,0.4) -- (2,4.55) node[above,font=\scriptsize,black] {світ змінився};
\draw[densely dashed,gray!60!black] (5,0.4) -- (5,4.55) node[above,font=\scriptsize,black] {перший сік за $B$};
\pic[scale=0.85] at (2,5.25) {nojuice};
\pic[scale=0.85] at (5,5.25) {juice};
% delta: dips after the change, burst at the first juice for B, then back to zero
\draw[thick,red!70!black] (0,3.65) -- (2.3,3.65) -- (2.3,3.3) -- (2.5,3.3) -- (2.5,3.65) -- (3.2,3.65) -- (3.2,3.3) -- (3.4,3.3) -- (3.4,3.65) -- (4.1,3.65) -- (4.1,3.35) -- (4.3,3.35) -- (4.3,3.65) -- (5.0,3.65) -- (5.0,4.35) -- (5.2,4.35) -- (5.2,3.65) -- (5.8,3.65) -- (5.8,4.1) -- (6.0,4.1) -- (6.0,3.65) -- (6.6,3.65) -- (6.6,3.85) -- (6.8,3.85) -- (6.8,3.65) -- (10,3.65);
% gain: high, drops after surprise, recovers
\draw[thick,blue!60!black] (0,3.2) -- (3.0,3.2) -- (3.4,2.75) -- (5.8,2.75) -- (7.2,3.2) -- (10,3.2);
% ACh: low, rises after the change, falls after learning
\draw[thick,green!45!black] (0,1.75) -- (2.8,1.75) -- (3.3,2.25) -- (6.8,2.25) -- (7.8,1.75) -- (10,1.75);
% choice of B
\draw[thick,black] (0,0.7) -- (3.0,0.7) plot[domain=3:10,samples=40] (\x,{0.7+0.6/(1+exp(-(\x-5.3)*1.8))});
\draw[->,gray!70!black] (0,0.2) -- (10.2,0.2) node[below left,font=\scriptsize,black] {час};
\end{tikzpicture}
\caption{Одна подія, простежена через усю машину. Це якісна схема за моделями й гіпотезами розділів~\ref{sec:dofa}--\ref{sec:acet}, а не результат розрахунку. Після зміни \nt{DOPA} відповідає провалами; повторні провали, за гіпотезою, піднімають сигнал несподіванки, \nt{NORA} знижує підсилення, \nt{ACET} вмикає запис; перший сік за $B$ дає сплеск, вибір переходить на $B$, і ручки повертаються.}
\label{fig:timeline}
\end{figure}

Зауважте, що жодна ручка не знає, що роблять інші. Кожна відгукується на свої ознаки --- помилку, її незвичну величину, невизначеність, --- і порядок дій складається сам, як в асинхронній схемі, де кожен блок спрацьовує, коли готові його входи. Таку злагоджену роботу загалом, наскільки нам відомо, ще не перевіряли: окремо ручки виміряно, а їхня спільна робота --- гіпотеза.

\section{Чим ця машина не схожа на комп'ютер}

\begin{table}[t]
\centering
\footnotesize
\setlength{\tabcolsep}{4pt}
\renewcommand{\arraystretch}{1.15}
\begin{tabular}{>{\raggedright}p{2.2cm}>{\raggedright}p{4.2cm}>{\raggedright\arraybackslash}p{6.6cm}}
\toprule
& Звичайний комп'ютер & Мозок \\
\midrule
час & спільний такт, близько гігагерца & спільного такту немає; кожен нейрон змінюється сам, вікна задають події й місцеві ритми (розд.~\ref{sec:glutcomputer}, \ref{sec:gaba}, \ref{sec:code}) \\
елементи & надійні двійкові вентилі & аналогові порогові елементи; синапс губить більшість імпульсів (розд.~\ref{sec:code}) \\
знак зв'язку & будь-який & задається нейроном-відправником (розд.~\ref{sec:neurontypes}) \\
пам'ять & окремо від процесора & у тих самих зв'язках, що й обчислення, і в самопідтримній активності (розд.~\ref{sec:glutcomputer}, \ref{sec:acet}) \\
навчання & зворотне поширення помилки & місцеві правила, одне широкомовне число $\delta$ (розд.~\ref{sec:dofa}) і проводи помилки до виходів кола рухів (розд.~\ref{sec:motorlearning}) \\
керування & регістри й шина команд & чотири широкомовні канали, кожен --- джгут із кількох ліній (розд.~\ref{sec:serocomputer}, \ref{sec:dofa}--\ref{sec:acet}) \\
енергія & десятки--сотні ватів на один процесор & близько $20$ ватів на весь мозок, у середньому близько імпульсу за секунду на нейрон (розд.~\ref{sec:code}) \\
\bottomrule
\end{tabular}
\caption{Мозок і звичайний комп'ютер.}
\label{tab:computer}
\end{table}

Таблиця~\ref{tab:computer} зводить головні відмінності. Кожна з них --- не вада, яку природа не встигла виправити, а частина одного рішення. Без спільного такту потрібні датчики «увімкнення---вимкнення» і детектори збігів. Ненадійним елементам потрібна надлишковість багатьох нейронів. Пам'яті, суміщеній з обчисленням, потрібен \nt{ACET}, щоб запис не псував читання. Навчанню без зворотних проводів потрібні \nt{DOPA} і шум. А енергетична межа в один імпульс за секунду змушує всю мову бути ощадливою й витрачати імпульси на новини.

\section{Чого ми не знаємо}

Найчесніше завершити це зведення переліком того, чого ми не знаємо. Кожен пункт --- задача для інженера.

\emph{Код.} Ми знаємо ємність імпульсного каналу і знаємо, що одержувач вибирає, яку частину повідомлення читати (розділ~\ref{sec:code}). Але наскільки кора користується точним часом імпульсів і чи є в нейронів «слова», невідомо.

\emph{Навчання через багато шарів.} Широкомовний сигнал навчає повільно, і що більше нейронів впливає на результат, то повільніше (твердження~\ref{prop:broadcastcost}). Як тоді кора налаштовує мільярди ваг у багатошарових колах, невідомо; є моделі, в яких помилку між шарами розносять пачки імпульсів~\cite{Payeur2021}. Можливо, частина відповіді --- в обчисленнях усередині гілок самого нейрона, де один нейрон поводиться як маленька багатошарова мережа: щоб повторити відповідь моделі одного нейрона кори, штучній мережі потрібно п'ять--вісім шарів~\cite{Beniaguev2021}, а гілки нейронів людини вміють обчислювати навіть виключне АБО~\cite{Gidon2020}. Інший кандидат --- самонавчання: якщо кора вчиться передбачати власні входи, помилка виникає на місці, в кожному шарі, і спільний сигнал для цього не потрібен (розділ~\ref{sec:dofa})~\cite{Keller2018}.

\emph{Що насправді передають широкомовні канали.} Для \nt{DOPA} є стандартна картина і шість її розширень (розділ~\ref{sec:dofaalt}); для \nt{NORA} і \nt{ACET} --- правдоподібні гіпотези (розділи~\ref{sec:nora}, \ref{sec:acet}); для \nt{SERO} --- здебільшого здогадки.

\emph{Узгодження в часі.} Ми бачили, як обчислити функцію, відміряти час від події й нарізати потік на вікна місцевим ритмом. Але як величезна мережа без спільного такту узгоджує роботу далеких ділянок, зрозуміло лише почасти.

\emph{Енергія.} Двадцять ватів на все. Ми розуміємо, чому код має бути розрідженим (твердження~\ref{prop:economy}), але збудувати машину, яка робила б те саме за такої самої потужності, поки що не вміє ніхто~\cite{MehonicKenyon2022}.

Мозок --- обчислювальна машина, зібрана не інженером, але за законами, які впізнає будь-який інженер: RC-кола, пороги, зворотні зв'язки, фільтри, шини й широкомовні регістри. Її схему ми прочитали лише частково. Дочитуватимуть її ті, хто вміє читати схеми.

\section{Що далі в книзі}

Цей розділ зібрав обчислювальне ядро машини. Решта розділів виходить за його межі. Розділи~\ref{sec:sensors} і~\ref{sec:recognition} --- про входи: як датчики перетворюють світло, звук і молекули на імпульси і як машина впізнає в них предмети. Розділи~\ref{sec:muscles}--\ref{sec:chemoutput} --- про виходи й те, що їх обслуговує: м'язи, біль як сигнал тривоги, навчання рухів окремими проводами помилки і повільний хімічний вихід у тіло. Розділ~\ref{sec:debugging} --- про те, як машину налагоджують ззовні, зокрема інтерфейсами мозок---комп'ютер. Розділи~\ref{sec:eetypes} і~\ref{sec:dendrites} ще раз сортують нейрони, тепер за електричною функцією і за кількістю входів. Розділ~\ref{sec:sleep} --- про те, що машина робить, коли відключена від світу, а розділи~\ref{sec:livingmachine} і~\ref{sec:dishbrain} --- про машини, зібрані з живих нейронів на чипі.

\section{Задачі}

Задачі цього розділу поєднують результати різних розділів: кожна спирається щонайменше на два з них.

\begin{exercise}[\lvlA]
\label{ex:10:1}
\emph{Шина і регістри.} Чотири широкомовні канали --- джгути приблизно по п'ять ліній (твердження~\ref{prop:bundle}); кожна лінія змінюється раз на секунду й розрізняє чотири рівні. Скільки років знадобилося б керуванню, щоб передати те, що шина \nt{GLUT} ($8.6\cdot10^{10}$ нейронів із частотою~\eqref{eq:energy}, точність $1\ms$ за~\eqref{eq:spikeentropy}) передає за секунду?
\end{exercise}

\begin{exercise}[\lvlB]
\label{ex:10:2}
\emph{Дві ручки на одній петлі.} У петлі~\eqref{eq:ei} ручки~\eqref{eq:machine} стоять на збудженні: $\tau_E\dot E=-E+g[(1-a)w_{EE}E-w_{EI}I+h]$, \nt{GABA} ними не керується. За чисел рис.~\ref{fig:ei} пачка \nt{NORA} піднімає $g$ до $6$. За якого найменшого $a$ петля стійка? Чи можна крутити ручки \nt{NORA} і \nt{ACET} незалежно?
\end{exercise}

\begin{exercise}[\lvlB]
\label{ex:10:3}
\emph{Звідки береться ціна широкомовлення.} Виходи $N$ нейронів $y_k=f(u_k)+\xi_k$ з незалежним гауссовим шумом дисперсії $\sigma^2$; $R\approx R_0+\sum_kR'_k\xi_k$ з однаковими $R'_k$, \nt{DOPA} розсилає $\delta=R-R_0$. Знайдіть відношення сигнал/шум поправки $\delta\,\xi_jx_j$ за одну спробу. Як із $N$ зростає кількість спроб?
\end{exercise}

\begin{exercise}[\lvlB]
\label{ex:10:4}
\emph{Зв'язати звук і спалах.} Звук приходить на шину через $5\ms$, спалах --- через $30\ms$ плюс затримка першого імпульсу~\eqref{eq:latency} ($\tau=10\ms$, $U$ від $1.2\theta$ до $4\theta$); фон на кожному вході --- $5$ імпульсів за секунду. Виберіть затримку звукової лінії й найменше вікно збігів для всіх контрастів. Скільки хибних зв'язувань за секунду дасть фон?
\end{exercise}

\begin{exercise}[\lvlB]
\label{ex:10:5}
\emph{Моделювання: хто помічає зміну.} Критик бачить $y=s+\text{шум}$ ($R=1$); з імовірністю $1/200$ $s$ стрибком набуває нового значення з дисперсією $J^2=4$. Змоделюйте фільтр Калмана з $q=0$, що піднімає $P$ до $J^2$, коли $E\leftarrow0.9E+\delta^2/(P+R)$ перевищує $20$. На скільки його помилка менша, ніж за найкращого сталого $\alpha$?
\end{exercise}

\begin{exercise}[\lvlA]
\label{ex:10:6}
\emph{Скільки спроб на зміну звички.} Мережа вивчила $Q_A=0.8$, $Q_B=0.2$ і вибирає за~\eqref{eq:softmax}, $\sigma=1$, $g=8$. Тепер $A$ дає сік з імовірністю $0.2$; $Q_A\leftarrow Q_A+\alpha(r-Q_A)$, $Q_B$ не змінюється. Через скільки виборів $A$ імовірність вибрати $A$ впаде до $0.6$ при $\alpha=0.1$ і $0.5$? А якщо \nt{NORA} знизить $g$ до $2$?
\end{exercise}

\begin{exercise}[\lvlC]
\label{ex:10:7}
\emph{Джгут замість проводу.} Виходи $y_k=w_k+\xi_k$, винагорода $R=-\sum_ky_k^2$; лінія джгута розсилає групі з $G$ нейронів помилку її виходів, $w_k\leftarrow w_k+\eta\,\delta_l\xi_k$. Покажіть, що за найкращого $\eta$ до $\sum w_k^2=\varepsilon\sum w_k^2(0)$ потрібно $\approx(G+2)\ln(1/\varepsilon)$ спроб. Скільки за $\varepsilon=0.01$ вчиться коло з $10^6$ нейронів на п'яти лініях, якщо спроба --- раз на три секунди?
\end{exercise}
